\section{Konzeptphase}
\label{sec:Konzeptphase}
In dieser Phase soll, vorbereitend für die Realisierungsphase, ein Konzept zur Umsetzung der bereits beschrieben Anforderungen
aufgezeigt werden. Inklusive jener aus \VerweisVL{sub:Zukünftig mögliche Anforderungen}. Dies beinhaltet ein technisches Konzept,
ein \Abkuerzung{UI}-Entwurf sowie die Aufstellung eines Testplans.


\subsection{Technisches Konzept}
\label{sub:Technisches Konzept}
Die Anforderungen werden in technische Arbeitspakete gruppiert, wie in \VerweisVL{sub:Vorgehensweise} erläutert. Daraus
resultieren die folgenden Arbeitspakete: Konfiguration von Atmosphere, Kommunikation: Server zu Client, Datenhaltung,
Datenübertragung, Event-Handling und \Abkuerzung{UI}-Entwurfsumsetzung. Im Folgenden werden einige Pakete auszugsweise genauer
spezifiziert.

\begin{description}
	\item [Kommunikation: Server zu Client]
Das Framework \Fachbegriff{Atmosphere} bietet weitreichende Features an, daher sollen bei der Umsetzung möglichst viele Aufgaben
vom Framework übernommen werden. Wie \zB der Verbindungsaufbau nach einem unerwarteten Abbruch und ein Nachrichten-Cache, um
Nachrichten neu zu versenden, die ein angemeldeter Client nicht erhalten hat. Ein angemeldeter Client bedeutet in diesem Kontext,
ein Browser der eine bestehende Verbindung zum Server hat. Auch das Senden von \sog "`Heartbeats"' (zu Deutsch: Herzschläge) kann
konfiguriert werden. Diese werden benötigt um die Verbindung bei zwischengeschalteten Proxies, welche lang suspendierte
Verbindungen unterbrechen, aufrecht zu erhalten.
	\item [Datenhaltung]
In einem Repository \VerweisKT{subsub:UserStatusRepository} sollen Benutzer-abhängige Informationen mit niedriger Kritikalität
gehalten werden. Eine niedrige Kritikalität liegt vor, wenn es akzeptabel ist, die Information \zB durch einen Serverneustart zu
verlieren. Diese wurde zuvor in der Datenbank abgelegt.
	\item [Datenübertragung]
Daten vom Server zum Client/Browser werden im Format \Abkuerzung{JSON}\footnote{Erklärung siehe \citep{JSON}.} übermittelt. Dazu
soll jede Event-Klasse eine generalisierte Klasse erweitern. Diese bietet eine entsprechende Methode an, welche die Event-Daten im
\Abkuerzung{JSON}-Format zurückliefert. Damit wird sichergestellt, dass übertragende Daten immer in diesem Format vorliegen. Es ist den
abgeleiteten Klassen erlaubt, diese Methode zu überschreiben, um den Inhaltsumfang selbst zu definieren.
\end{description}

\subsection{UI-Entwurf}
\label{sub:UI-Entwurf}
Die grafische Benutzeroberfläche ist ein zentraler Bestandteil des \Abkuerzung{VMS}, weil sie die Schnittstelle zwischen Anwender
und Software ist. Da bereits eine Weboberfläche vorhanden ist, müssen Änderungen in diese integriert werden. Der
Gestaltungsspielraum wird dadurch stark eingeschränkt: Es ist daher verstärkt darauf zu achten, dass die Änderungen gut mit der
restlichen Oberfläche harmonisieren. Da diese größtenteils aus Graustufen besteht werden ausschließlich Pastelltöne verwendet.\\
Die Umsetzung der \Abkuerzung{UI} erfolgte prototypisch und ist im \Anhang{app:UI-Entwurf} zu sehen. Die Idee aus
\VerweisKL{sub:Soll-Konzept} dem Benutzer auf explizit neue Tasks hinzuweisen, wurde dabei bereits aufgenommen. Es wurde
vernachlässigt unterschiedliche Formen pro Task-Typ zu entwerfen.


\subsection{Proof-of-Concept}
\label{sub:Proof-of-Concept}
Es wurde eine Demo-Implementierung mit dem Framework Atmosphere vorgenommen. Dazu wurde ein Chat-Beispiel nachgebaut. Das
Augenmerk lag dabei auf der Einrichtung von Atmosphere, sowie der richtigen Konfiguration der Webapplikation für eine asynchrone
Verarbeitung. Weil keine Minderung der Serverperformance festgestellt wurde, konnte mit dem Framework fortgefahren werden.


\subsection{Testplan}
\label{sub:Testplan}
In einem Testplan werden die definierten Anforderungen zusammengefasst und beschrieben wie diese zu prüfen sind. Dies sichert
bereits vor der Realisierungsphase die vollständige Umsetzung aller Anforderungen. Neben den in der \TabelleK{tab:Testfälle}
aufgeführten Testfällen, soll mit einer sinnvollen Testabdeckung durch JUnit-Tests eine hohe Softwarequalität und -robustheit
erreicht werden.
\newTabelle{Testfälle}{tab:Testfälle}{Testfälle.tex}